%% ma: L10-B14-0048
%% lop: 10
%% chuong: 5
%% bai: 14
%% dang: 10.14.8
%% loai: TL
%% muc: VD
%% dap_an: "-"
%% nguon: "chuyen-de-cac-so-dac-trung-cua-mau-so-lieu-khong-ghep-nhom-toan-10-chuong-trinh-moi.pdf (GV. Phan Nhật Linh) (P05), Bài 03 Dạng 2 Bài tập tự luận 8"
%% kien_thuc: "Lớp 10 Bài 13 (số trung bình), Bài 14 (phương sai, độ lệch chuẩn)"
%% trang_thai: chinh-thuc
%% ghi_chu: "Lời giải gốc làm tròn trung bình thành 155 khi tính; kết quả 48,15 không đổi. Tự duyệt theo yêu cầu giáo viên 22/10/2026."
\begin{ex}
Một cửa hàng bán gạo thống kê số ki-lô-gam gạo bán được mỗi ngày trong $30$ ngày, được bảng tần số sau:
\begin{center}
\begin{tabular}{|c|c|c|c|c|c|c|c|c|}\hline
Số kg gạo & $100$ & $120$ & $130$ & $160$ & $180$ & $200$ & $250$ & Tổng\\\hline
Tần số & $7$ & $4$ & $2$ & $8$ & $3$ & $2$ & $4$ & $30$\\\hline
\end{tabular}
\end{center}
Tính độ lệch chuẩn của mẫu số liệu (làm tròn đến hàng phần trăm).
\loigiai{$\overline x=\dfrac{4660}{30}=\dfrac{466}3\approx155{,}33$.

$s^2=\dfrac1{30}\sum n_ix_i^2-\overline x^2=\dfrac{793\,400}{30}-\left(\dfrac{466}3\right)^2\approx2318{,}22$.

$s\approx48{,}15$.}
\end{ex}
